\section{Computational verification protocol}\label{app:computation}

The companion directory \texttt{anc/supplement/} contains five groups of Python scripts in the subdirectories \texttt{part1\_comparison}, \texttt{part2\_fixed\_kernel}, \texttt{part3\_two\_term} \texttt{part4\_optimality} and \texttt{part5\_revision}, their captured outputs, a README recording the interpreter and package versions, an MIT licence for the computational files, and a \texttt{SHA256SUMS} manifest covering every file except the manifest itself.  The top-level runner \texttt{run\_all\_checks.py} executes all fifteen distinct certificates in a temporary copy and compares their outputs with the captured files.  All arithmetic is exact, over \(\mathbb Q\) or a prime field; the two scripts \texttt{verify\_nakayama\_hom\_table.py} and \texttt{verify\_a63\_stable\_kernel.py} use SymPy for exact rational matrices, and every other script uses only the standard library.  The scripts are verification aids; the proofs in the text do not depend on executing them.

\subsection*{Comparison and Nakayama computations}
The file \texttt{verify\_nakayama\_hom\_table.py} implements the path-space calculation in \Cref{thm:nakayama-hom-table}: right modules over \(k\widetilde Q_{2d}/\mathfrak r^d\), cohomological cones \([P_{i+1}\to P_i]\) in degrees \(-1,0\), the middle mapping-complex terms ordered as in \eqref{eq:nakayama-mapping-complex}, and differentials \(\delta^{-1}(f)=(f\alpha_i,\alpha_jf)\), \(\delta^0(a,b)=\alpha_ja-b\alpha_i\).  For every offset and for \(d=3,4,5\) it constructs the two matrices over \(\mathbb Q\), verifies that their composite is zero, computes their ranks exactly, and compares the cohomology dimensions with the theorem; \texttt{verify\_nakayama\_hom\_table\_char2.py} repeats the checks over \(\mathbb F_2\).  The file \texttt{verify\_a63\_stable\_kernel.py} represents every uniserial module \(M(i,\ell)\) of \(\widetilde A^3_6\) as a cyclic quiver representation, solves the arrow-commutativity equations for Hom-spaces, computes the subspace of maps factoring through the six projectives, and verifies the stable Hom formula of \Cref{lem:a63-stable} for shifts \(-8\le q\le8\).  Three further files check the two module-side orientations of the AHKLY presentation (\texttt{verify\_ahkly\_sides.py}), the one-object duality maps of \eqref{eq:one-object-obs} (\texttt{verify\_perfect\_complex\_duality.py}), and the factorisation criterion on small examples (\texttt{verify\_factorisation\_toys.py}).

\subsection*{Fixed-kernel computations}
\texttt{check\_7dim.py}: for \(C=k[x,y]/(x^2,y^2)\), the kernels of \(x\) and \(y\), the Tor groups of \Cref{ex:seven}, the stable endomorphism ring of \(C/(y)\), and the infinite projective dimension of \(k_A\); its variable \texttt{N} denotes the test object \(X=C/(y)\).
\texttt{verify\_dg\_boundary.py}: the totalised semi-free resolutions of \(L_1\) and \(L_4\) over the dg algebra of \Cref{thm:dg-boundary}, the check \(d^2=0\), the cohomology, the nonzero degrees of \(\Ext^q_B(L_1,L_1)\) and \(\Ext^q_B(L_4,L_1)\) for \(0\le q\le32\), and an independent stable-Ext computation over \(A_3\).
\texttt{verify\_nakayama\_profiles.py}: all stable Hom spaces between the uniserial modules of \(A_3\) and \(A_4\) over \(\mathbb Q\), the check of \Cref{lem:stable-hom}(ii) on \(144\) and \(576\) entries, and every realising module of \Cref{prop:m3} over a full shift period.
\texttt{verify\_formal\_realisation.py}: for five pairs \((\Gam,U)\), including a noncommutative \(\Gam\) and a nonprojective bimodule, associativity of the ring \eqref{eq:R-delta} with \(\delta=0\), multiplicativity of \(\lambda\), the two module decompositions, and the strict dg map on products.
\texttt{verify\_group\_block\_obstruction.py}: over \(\mathrm{GF}(3)\), the block decomposition of \(k[C_3\times C_2]\), the periodic resolutions of the two simples, and the vanishing of all cross-Ext groups used after \Cref{thm:pgroup}.
The runner \texttt{run\_checks.py} re-executes these scripts and compares each output byte for byte with the captured file.

\subsection*{Two-term kernel computations}
\texttt{verify\_nonformal.py}, with the helpers \texttt{exact.py} and \texttt{m3.py}: for \Cref{ex:nonformal}, the associativity, bimodule, Peiffer and Leibniz identities of the nine-dimensional dg algebra; the multiplicativity of \(\lambda_\delta\) on all basis pairs; the decompositions \(S_{R}\cong(pR)^{\oplus2}\) and \({}_RS\cong(Re)^{\oplus2}\) on all basis actions; the strict dg map \(D\to\End^\bullet_R(T)^{\op}\) on all products and differentials; and the homotopy transfer of \(m_3\), computed once on \(D\) and once on the presentation-cone algebra, together with the rank of the Hochschild coboundary matrix, which shows that the value \(w\) is not a coboundary.  The captured output is included.


\subsection*{Optimality and Schur-boundary computations}
The additional certificates cover the following finite calculations.  The file \texttt{verify\_connective\_obstruction.py} checks the Frobenius pairings of \Cref{ex:negative-cy-algebras}, the complete endomorphism complex, strict dg map and vertex-component cohomology of \Cref{ex:sphere-schur}, and the analogous sixteen-dimensional tensor Schur example for the exterior algebra.  It also checks the shifted two-term Hom benchmark and an eighteen-dimensional non-formal tensor control, including its Massey-product indeterminacy.  The helpers \texttt{exact.py}, \texttt{m3.py} and \texttt{verify\_nonformal.py} are copies of those in \texttt{part3\_two\_term}.  Complete basis multiplication tables for both Schur rings are supplied as JSON files.  All degrees of every finite complex are checked, with independent rank checks over \(\mathbb F_2,\mathbb F_3,\mathbb F_{101}\); the arbitrary-field arguments are the proofs in the text.  The file \texttt{verify\_marking\_boundary.py} checks the full Hom complex and multiplication of the connective marking in \Cref{ex:marking-boundary}.  Neither script purports to test all ambient rings or to infer homologicality from finitely many Tor groups.  The nonexistence assertions in \Cref{thm:texact-optimality,prop:negative-cy-obstruction} are proved mathematically.

\subsection*{Coordinate symmetry}
The file \texttt{verify\_coordinate\_swap.py} checks the unital algebra involution of \Cref{ex:support-symmetry}, the exchanged module actions, and the image--kernel identities for one repeating step of the free and Hom resolutions over \(\mathbb Q\) and \(\mathbb F_2\). The proof of the infinite resolution and of the induced equivalence of marked cuts is in the text; these finite checks do not replace it.
